\documentclass{article}
\usepackage[T1]{fontenc}
\usepackage{lmodern}
\usepackage{graphicx}
\usepackage{amsmath,amssymb}
\usepackage[a4paper,margin=2cm]{geometry}
\usepackage[absolute,overlay]{textpos}
\setlength{\TPHorizModule}{1mm}
\setlength{\TPVertModule}{1mm}
\usepackage[indonesian]{babel}
\usepackage{hyperref}
\setlength{\emergencystretch}{2em}
% unit: Halaman 55

\begin{document}

% === Halaman 55 (python/native) ===

Operasi aritmetika pada GF(2m) 

 Misalkan a = (am-1..a1 a0) dan b = (bm-1...b1 b0) GF(2m) 

•

Penjumlahan:  


a + b = c = (cm-1..c1 c0) dimana ci = (ai + bi) mod 2, c  GF(2m) 

•

Perkalian: a  b = c = (cm-1..c1 c0 ) dimana c adalah sisa pembagian polinom 

a(x)  b(x) dengan irreducible polynomial derajat m,  c  GF(2m)

•

Algoritma Rijndael (AES) beroperasi di dalam GF(28)

\end{document}
